%=======================================================================
%  IE 5311-001 (D01) Exploration Problem EP7
%  One page, single spacing, 10 pt, 1-inch margins, as the sheet requires.
%  Content past page 1 is not graded, so the page limit is hard.
%  Sections map onto the three parts asked: algorithm, derived
%  complexity, and when DP beats a general-purpose IP solver.
%
%  IF IT SPILLS ONTO A SECOND PAGE, delete in this order until it fits.
%    1. the last sentence of section 3 (from "A solver also returns")
%    2. section 4 entirely (verification is beyond what was asked)
%    3. the $10^2$ and $10^3$ rows of the scaling table
%  Do not cut section 2 or section 3: they are two of the three things
%  the question explicitly asks for.
%=======================================================================
\documentclass[10pt]{article}

\usepackage[margin=1in]{geometry}
\usepackage{amsmath}
\usepackage{amssymb}
\usepackage{booktabs}
\usepackage{array}

\pagestyle{empty}
\setlength{\parskip}{0.35em}
\setlength{\parindent}{0pt}

\begin{document}

{\large\bfseries EP7.\quad Dynamic programming for the knapsack problem}\\[1pt]
{\small Lidivine Kengne \quad\textbar\quad IE 5311-001 (D01), Homework 5}
\vspace{0.3em}

\textbf{1. The algorithm.}
Items $i \in \{1,\dots,n\}$ carry values $v_i$ and integer weights
$w_i$; the capacity $W$ is an integer. The 0--1 knapsack problem is
$\max \sum_i v_i y_i$ subject to $\sum_i w_i y_i \le W$,
$y \in \{0,1\}^n$. Let $f(i,c)$ be the best value obtainable from the
first $i$ items under capacity $c$. Then $f(0,c) = 0$ and
\[
  f(i,c) \;=\;
  \begin{cases}
    f(i-1,c), & w_i > c,\\[2pt]
    \max\bigl\{\, f(i-1,c),\; v_i + f(i-1,\,c-w_i) \,\bigr\}, & w_i \le c,
  \end{cases}
\]
the branches being \emph{leave} and \emph{take}. It is valid because a
partial solution affects the rest only through the capacity it leaves
free, so $(i,c)$ is a sufficient state. On the Problem~2 pricing
instance --- $a = (2,2,5,6,2,5,5)$, $b = 12$, values
$\pi_i - c_{i1} = (21.7,-14.0,-1.3,-5.7,-0.3,7.3,7.7)$ --- the table is
$8 \times 13$ and its rows $2$ to $5$ reproduce row $1$ exactly,
because those four items carry negative value, so \emph{leave} always
wins and the table shows directly that no optimum takes them.
Backtracking reads the solution off it: item $i$ is taken exactly when
$f(i,c) \neq f(i-1,c)$, giving $y = (1,0,0,0,0,1,1)$ of weight
$2+5+5 = 12$, filling the knapsack exactly at value
$f(7,12) = 36.7$.

\textbf{2. Complexity, derived.}
The state space is $\{0,\dots,n\} \times \{0,\dots,W\}$, of size
$(n+1)(W+1)$. Each state is evaluated once, and its evaluation is one
comparison between two entries already computed, hence $\Theta(1)$
work. Total time is therefore $\Theta(nW)$ and storage $\Theta(nW)$,
the latter reducible to $\Theta(W)$ by keeping one row and sweeping $c$
downward so item $i$ is not reused; backtracking adds $\Theta(n)$.
Compare that with the \emph{input length}: encoding the instance takes
$\Theta\bigl(n(\log v_{\max} + \log w_{\max}) + \log W\bigr)$ bits, so
$W$ contributes $\log W$ bits but a factor $W$ to the running time.
The algorithm is thus exponential in the size of its input,
$\Theta(nW) = \Theta(n2^{\log W})$, which is what
\emph{pseudo-polynomial} names, and it contradicts nothing: knapsack
is NP-hard. Against the $\Theta(2^n)$ of enumeration the gain is real
all the same --- above, $104$ cells replace $2^7 = 128$ subsets.

\textbf{3. When to prefer it to an IP solver.}
Four conditions favour the table. (i) Weights are integers and $W$ is
modest, so $nW$ stays small. (ii) The same structure is solved
repeatedly with only the objective changing, so the DP rebuilds
nothing and runs a fixed number of operations where a solver pays
model construction, presolve and a search on every call. This is
exactly the Problem~2 pricing subproblem: for agent $j$ the least
reduced cost column solves
$\min_{y} \sum_i (c_{ij} - \pi_i) y_i - \mu_j$ subject to
$\sum_i a_i y_i \le b$, that is
$-\max_{y} \sum_i (\pi_i - c_{ij}) y_i - \mu_j$, a knapsack whose
weights $a_i$ and capacity $b$ never change while the values
$\pi_i - c_{ij}$ change every iteration. Negative values need no
special treatment, since \emph{leave} is always available.
(iii) Integer values make the recursion exact in integer arithmetic,
with no tolerance and no risk of a wrong incumbent. (iv) Runtime is
deterministic and known in advance, where a search tree is not.
Scaling the weights and capacity by a common factor leaves the
solution unchanged but grows $W$, which prices (i) directly
(\texttt{ep7\_knapsack\_dp.py}):

\begin{center}\footnotesize
\setlength{\tabcolsep}{7pt}
\begin{tabular}{@{}rrrrr@{}}
\toprule
scale & $W$ & cells & DP (ms) & Gurobi (ms)\\
\midrule
$1$      &      12 &       104 & 0.03 & 0.22\\
$10$     &     120 &       968 & 0.03 & 0.19\\
$10^{2}$ &  1\,200 &   9\,608 & 0.03 & 0.20\\
$10^{3}$ & 12\,000 &  96\,008 & 0.15 & 0.21\\
$10^{4}$ & 120\,000 & 960\,008 & 2.69 & 0.25\\
\bottomrule
\end{tabular}
\end{center}

The DP is flat while the table fits in cache, then linear in $W$ as
$\Theta(nW)$ predicts, while the solver barely moves; the crossover is
near $W \approx 2 \times 10^{4}$ for $n = 7$. A solver also wins once
weights are fractional, since the table does not exist, and once side
constraints appear, since extra capacities multiply the state space
while a solver absorbs extra rows.

\textbf{4. Verification.}
\texttt{ep7\_knapsack\_dp.py} substitutes the recursion for the Gurobi
pricer in \texttt{p2\_colgen.py} and re-runs column generation. Over
the 39 pricing calls made, the dynamic program, the Gurobi binary
program and enumeration of all 128 subsets returned the same reduced
cost every time, and the run still reached $61.000000$. Pricing cost
$0.041$\,ms per call by table against $0.549$\,ms by solver, a factor
of $13.5$: the DP vectorised over capacities, the solver timed as
column generation actually calls it, model construction included.

\end{document}
